% !TEX program = xelatex
\documentclass[12pt,a4paper,landscape]{article}

\usepackage[UTF8,heading=true]{ctex}
\setCJKmainfont{SimSun}[BoldFont=SimHei,ItalicFont=KaiTi]
\setCJKsansfont{SimHei}

\usepackage[top=20mm,bottom=20mm,left=20mm,right=20mm,headheight=14pt,landscape]{geometry}
\usepackage{tikz}
\usetikzlibrary{arrows.meta,positioning,calc,shapes.geometric}
\usepackage{fontawesome5}
\usepackage{multicol}
\usepackage{amsmath}
\usepackage{enumitem}
\usepackage{booktabs}
\usepackage{listings}
\usepackage{xcolor}
\usepackage{hyperref}
\usepackage{fancyhdr}
\usepackage{titlesec}
\usepackage{tcolorbox}

\tcbuselibrary{skins,breakable}

\hypersetup{colorlinks=true, linkcolor=blue!60!black, urlcolor=blue!60!black}

% 代码样式
\lstset{
  basicstyle=\ttfamily\small,
  backgroundcolor=\color{gray!8},
  frame=single,
  rulecolor=\color{gray!30},
  breaklines=true,
  numbers=none,
  xleftmargin=4mm,
  framexleftmargin=2mm,
}

% 页眉
\pagestyle{fancy}
\fancyhf{}
\fancyhead[L]{\small\color{gray} 动力装备知识图谱RAG系统 — 第一阶段汇报}
\fancyhead[R]{\small\color{gray} 纪文龙}
\fancyfoot[C]{\thepage}

% 术语解释框
\newtcolorbox{termbox}[1]{
  colback=blue!3, colframe=blue!40, fonttitle=\bfseries,
  title={#1}, rounded corners, boxrule=0.5pt,
  left=4pt, right=4pt, top=2pt, bottom=2pt,
  before skip=4pt, after skip=4pt,
}

\begin{document}

% ============================================================
\begin{center}
  {\LARGE\bfseries 第一阶段工作详解 — 看完你就能汇报}\\[8pt]
  {\large 纪文龙 \quad 2026年4月10日}\\[4pt]
  {\small GitLab分支: \texttt{jiwenlong/vectorize-pipeline} \quad 提交: \texttt{2a2b473}}
\end{center}
\vspace{4mm}

\noindent\rule{\textwidth}{0.4pt}

\textbf{本文档用途}：用大白话帮你搞懂你「做了什么」和「组会怎么讲」。每个专业术语都有解释框，看不懂就看蓝框。

\noindent\rule{\textwidth}{0.4pt}
\vspace{4mm}

% ============================================================
\part*{第一部分：你到底做了什么}
% ============================================================

\section{一句话总结}

你做了一个\textbf{自动化处理流程}，效果是：给它一篇论文的OCR标注数据，它能自动存起来，然后你问它问题，它能找到论文中最相关的段落回答你。

\begin{termbox}{什么是 Pipeline（管线/流水线）？}
Pipeline 就是「一步一步处理的流水线」，像工厂的生产线一样：原材料进去 → 经过几道工序 → 出成品。在我们这里：\\[2pt]
\quad 原材料 = 标注好的JSON文件（论文的OCR识别结果）\\
\quad 成品 = 一个能回答问题的系统\\[2pt]
\textbf{汇报时你可以说}：「我做了一个数据处理的流水线」，不用说Pipeline这个词也完全可以。
\end{termbox}

\section{整体流程图}

下面这张图就是你的「流水线」，从左到右依次经过6个步骤：

\begin{center}
\begin{tikzpicture}
  \tikzset{sbox/.style={rectangle, rounded corners=6pt,
    minimum width=28mm, minimum height=18mm, align=center, font=\small}}

  \node[sbox, draw=blue!70, fill=blue!7] (s1) at (0,0) {\faFileCode\\[2pt] \textbf{解析JSON}\\[-1pt]{\scriptsize 246块}};
  \node[sbox, draw=orange!70, fill=orange!7] (s15) at (3.5,0) {\faBroom\\[2pt] \textbf{数据清洗}\\[-1pt]{\scriptsize 228块}};
  \node[sbox, draw=purple!70, fill=purple!7] (s2) at (7,0) {\faCut\\[2pt] \textbf{文本分块}\\[-1pt]{\scriptsize 16 chunk}};
  \node[sbox, draw=green!50!black, fill=green!7] (s3) at (10.5,0) {\faProjectDiagram\\[2pt] \textbf{BGE-m3}\\[-1pt]{\scriptsize 变成数字}};
  \node[sbox, draw=blue!40!black, fill=blue!5] (s4) at (14,0) {\faDatabase\\[2pt] \textbf{Chroma}\\[-1pt]{\scriptsize 存起来}};

  \draw[-{Stealth[length=4pt]},blue!70,line width=1.5pt] (s1) -- (s15);
  \draw[-{Stealth[length=4pt]},orange!70,line width=1.5pt] (s15) -- (s2);
  \draw[-{Stealth[length=4pt]},purple!70,line width=1.5pt] (s2) -- (s3);
  \draw[-{Stealth[length=4pt]},green!50!black,line width=1.5pt] (s3) -- (s4);

  \node[sbox, draw=red!70, fill=red!7] (s5) at (5,-2.8) {\faSearch\\[2pt] \textbf{混合检索}\\[-1pt]{\scriptsize 找相关段落}};
  \node[sbox, draw=orange!80, fill=orange!7] (s6) at (9.5,-2.8) {\faComments\\[2pt] \textbf{RAG问答}\\[-1pt]{\scriptsize 回答问题}};

  \draw[-{Stealth[length=4pt]},blue!40!black,line width=1.2pt] (s4.south) -- ++(0,-0.8) -| (s5.north);
  \draw[-{Stealth[length=4pt]},red!70,line width=1.2pt] (s5) -- (s6);

  \node[rectangle, rounded corners=3pt, draw=gray, font=\small\color{gray}] (q) at (1.5,-2.8) {\faUser\; 用户问问题};
  \draw[-{Stealth[length=3pt]}, gray, dashed] (q) -- (s5);
\end{tikzpicture}
\end{center}

每一步是什么意思？下面逐个解释。先说一个贯穿全文的词：

\begin{termbox}{什么是「文本块」？}
一篇论文的PDF扫描件上有很多段文字。郭继鸿用标注工具在每一段文字上画了一个框，每个框里的文字就是一个\textbf{文本块}。\\[2pt]
比如一个标题是一个块，一个段落是一个块，一个列表项也是一个块。4页的论文一共画了246个框，所以就有246个文本块。\\[2pt]
你可以把「文本块」理解为「一小段文字」。
\end{termbox}

\begin{termbox}{什么是 OCR？}
OCR = 光学字符识别。就是把纸质文件或PDF扫描件上的文字，用软件自动识别成电脑能编辑的文字。\\[2pt]
我们项目里的论文都是扫描版PDF（相当于照片），电脑看不懂里面的文字，需要OCR先把文字「读」出来，郭继鸿用Label Studio做的就是这个事情。
\end{termbox}

% ============================================================
\section{第1步：解析JSON}

\begin{termbox}{什么是 Label Studio JSON？}
郭继鸿用一个叫 Label Studio 的标注工具，在PDF论文的扫描图片上画框，标记出「这段是标题」「这段是正文」「这段是列表」。标记完成后导出一个 JSON 文件（一种数据格式，类似于Excel表格但更灵活）。\\[2pt]
这个JSON文件就是我们的\textbf{原材料}。
\end{termbox}

\textbf{这一步做了什么}：读取JSON文件，把里面的文字提取出来。每一条标注包含：
\begin{itemize}[nosep]
  \item 位置信息（在页面上的坐标）
  \item 类型标签（Title=标题 / Para=段落 / List=列表项）
  \item 文字内容（OCR识别出的文字）
\end{itemize}

\textbf{一个改进}：原版代码只处理了「标题」和「段落」两种标签，\textbf{漏掉了「列表」}（就是论文里的编号列表，比如 1. 2. 3. 这种），有15个列表段落被跳过了。我的代码补上了。

\textbf{结果}：从1篇论文（4页）中提取出246个文本块（= 246小段文字）。

% ============================================================
\section{第1.5步：数据清洗}

\begin{termbox}{什么是数据清洗？}
原始数据里有「脏数据」——不完整、有错误、质量差的数据。数据清洗就是把这些脏数据修好或删掉，就像洗菜一样，把烂叶子摘掉。
\end{termbox}

\textbf{发现了什么问题？}

因为这篇论文是\textbf{双栏排版}（左右两列文字），OCR在画框时把一些句子切断了，产生了碎片：

\begin{center}
\begin{tabular}{lll}
\toprule
\textbf{碎片} & \textbf{页码} & \textbf{怎么回事} \\
\midrule
``程体系。'' & 第1页 & 「课程体系。」被切成两半，只留了``程体系。'' \\
``建方法''   & 第1页 & 「构建方法」被切成两半 \\
``印证。''   & 第2页 & 句子末尾的碎片 \\
\bottomrule
\end{tabular}
\end{center}

\textbf{怎么修的？}

我写了一个函数：如果一个文本块不足10个字，就自动合并到旁边的块里。

\textbf{效果}：246块 $\to$ 228块，合并了18个碎片。之后质量检测显示 0 个问题。

% ============================================================
\section{第2步：文本分块}

\begin{termbox}{什么是分块（Chunking）？}
228个小块太碎了，直接存不好检索。分块就是把很多小块拼成合适大小的「段落」——大约500个字符一段（约250个中文字），既不太长也不太短。拼好的段落叫 \textbf{chunk}。\\[2pt]
\textbf{为什么要这个大小？}太短的话信息不够，大模型理解不了上下文；太长的话搜索时噪音太多，无关内容太多。500字符是业界常用的经验值。
\end{termbox}

\textbf{一个关键改进：标题触发新段落}

如果不做特殊处理，可能出问题：

\begin{center}
\begin{tikzpicture}[font=\small]
  \node[font=\normalsize\bfseries\color{red!70!black}] at (-5.5,2) {\faTimesCircle\; 普通做法（有问题）};
  \node[rectangle, draw=red!40, fill=red!5, rounded corners=3pt,
    text width=60mm, inner sep=5pt, align=left] at (-5.5,0.5) {
    {\color{gray} ...前面的正文正文正文...}\\[2pt]
    {\color{red!70!black}\bfseries 一、新章节的标题}\\[1pt]
    {\color{gray}\scriptsize $\uparrow$ 标题粘在上一段末尾了！}
  };
  \node[rectangle, draw=gray!30, fill=gray!5, rounded corners=3pt,
    text width=60mm, inner sep=5pt, align=left] at (-5.5,-1.5) {
    {\color{gray} 新章节的正文...}\\[1pt]
    {\color{gray}\scriptsize $\uparrow$ 这段没有标题了！}
  };

  \node[font=\normalsize\bfseries\color{green!60!black}] at (5.5,2) {\faCheckCircle\; 我的做法（正确）};
  \node[rectangle, draw=gray!30, fill=gray!5, rounded corners=3pt,
    text width=60mm, inner sep=5pt, align=left] at (5.5,0.5) {
    {\color{gray} ...前面的正文正文正文...}\\[1pt]
    {\color{gray}\scriptsize $\uparrow$ 上一段干净地结束}
  };
  \node[rectangle, draw=green!40, fill=green!5, rounded corners=3pt,
    text width=60mm, inner sep=5pt, align=left] at (5.5,-1.5) {
    {\color{green!60!black}\bfseries 一、新章节的标题}\\[2pt]
    {\color{gray} 新章节的正文...}\\[1pt]
    {\color{gray}\scriptsize $\uparrow$ 标题和正文在一起}
  };

  \node[font=\Large\bfseries\color{gray}] at (0,0) {VS};
\end{tikzpicture}
\end{center}

\textbf{为什么这很重要？}如果有人搜「新章节的标题」，左边的做法会匹配到上一段（因为标题在那里），但上一段的内容跟这个标题没关系。右边的做法就不会搞混了。

\textbf{结果}：228个小块拼成了16个chunk，平均长度429字符。

% ============================================================
\section{第3步：向量化}

\begin{termbox}{什么是向量化（Embedding）？}
电脑不认识中文，只认识数字。向量化就是把一段文字转成一串数字（叫做「向量」），这串数字能代表这段文字的\textbf{语义}（含义）。\\[4pt]
意思相近的文字，转出来的数字也很接近。比如：\\
\quad「燃气轮机状态监控」和「发动机健康管理」→ 数字很接近\\
\quad「猫吃鱼」和「燃气轮机」→ 数字差很远\\[4pt]
\textbf{用什么模型做的？}用的是 \textbf{BGE-m3}，北京智源研究院做的开源模型，目前中文检索效果最好的模型之一。每段文字会变成1024个数字（1024维向量）。
\end{termbox}

\begin{termbox}{什么是 Chroma？}
Chroma 是一个\textbf{向量数据库}——专门用来存那些数字（向量）的数据库。存进去之后，你给它一个新的问题，它能快速找到数据库里跟这个问题最接近的几段文字。\\[2pt]
类比：普通数据库像图书馆的编号检索（按书名、作者找书）；向量数据库像一个特别聪明的图书管理员（按「你想看什么内容的书」找书）。
\end{termbox}

\textbf{这一步做了什么}：
\begin{enumerate}[nosep]
  \item 加载BGE-m3模型
  \item 把16个chunk都转成1024维的数字（向量化）
  \item 存入Chroma数据库（保存在本地磁盘，下次不用重新算）
\end{enumerate}

% ============================================================
\section{第4步：混合检索}

这是最重要的升级。

\begin{termbox}{什么是检索（Retrieval）？}
检索就是「搜索」。你输入一个问题，系统在数据库里找到跟这个问题最相关的几段文字返回给你。跟百度搜索类似，但搜的是我们自己存进去的论文内容。
\end{termbox}

基础版本只用了一种搜索方法（语义检索）。我的版本用了\textbf{两种方法}，然后把结果合并：

\begin{center}
\begin{tikzpicture}[font=\small]
  \node[rectangle, rounded corners=6pt, draw=blue!60, fill=blue!5,
    minimum width=70mm, minimum height=28mm, align=center]
    (sem) at (0,0) {\textbf{方法一：语义检索}\\[3pt]把问题也变成1024个数字\\然后找数据库里数字最接近的段落\\[2pt]{\color{blue!60}\small 优点：理解同义词（发动机$\approx$轮机）}\\{\color{red!60}\small 缺点：精确关键词可能找不到}};
  \node[rectangle, rounded corners=6pt, draw=orange!60, fill=orange!5,
    minimum width=70mm, minimum height=28mm, align=center]
    (bm) at (9,0) {\textbf{方法二：BM25关键词检索}\\[3pt]把问题切成关键词\\然后找含有这些关键词的段落\\[2pt]{\color{blue!60}\small 优点：精确匹配（搜``布鲁姆''就找``布鲁姆''）}\\{\color{red!60}\small 缺点：不理解同义词}};

  \node[rectangle, rounded corners=6pt, draw=green!50!black, fill=green!5,
    minimum width=60mm, minimum height=18mm, align=center]
    (rrf) at (4.5,-3.5) {\textbf{合并：RRF融合排序}\\[2pt]70\%听语义检索的 + 30\%听BM25的\\两种方法互补，覆盖盲区};

  \draw[-{Stealth[length=5pt]},blue!60,line width=1.5pt] (sem.south) -- (rrf.north west);
  \draw[-{Stealth[length=5pt]},orange!60,line width=1.5pt] (bm.south) -- (rrf.north east);
\end{tikzpicture}
\end{center}

\begin{termbox}{什么是 BM25？什么是 jieba？什么是 RRF？}
\textbf{BM25}：一个经典的搜索算法，跟百度早期的搜索原理类似。你搜``燃气轮机''，它就去数几篇文章里``燃气轮机''出现了几次，出现多的排前面。\\[2pt]
\textbf{jieba}（结巴分词）：中文没有空格，电脑不知道从哪断词。jieba就是一个切词工具，比如把"数据库系统"切成"数据库/系统"。BM25需要先切词才能搜索。\\[2pt]
\textbf{RRF}（Reciprocal Rank Fusion）：一种合并排名的方法。语义检索给出一个排名、BM25给出一个排名，RRF把两个排名合在一起，算出最终排名。简单理解：两种方法投票，得票多的排前面。
\end{termbox}

\textbf{效果对比}——为什么两种方法都需要：

\begin{center}
\begin{tabular}{lccl}
\toprule
\textbf{搜索什么} & \textbf{语义检索得分} & \textbf{BM25得分} & \textbf{说明} \\
\midrule
``燃气轮机健康管理'' & 0.685（高） & 3.2 & 两种都能找到 \\
``布鲁姆认知过程'' & 0.387（低） & \textbf{4.7}（高） & 语义找不到，BM25能！ \\
``船舶动力装置'' & \textbf{0.542}（高） & 0.0（没有） & BM25找不到，语义能！ \\
\bottomrule
\end{tabular}
\end{center}

\textbf{汇报时你可以这样说}：「两种方法各有优劣，合在一起可以覆盖各自的盲区。」

% ============================================================
\section{第5步：RAG问答}

\begin{termbox}{什么是 RAG？}
RAG = Retrieval-Augmented Generation，翻译成「检索增强生成」。\\[2pt]
简单说：你问一个问题 → 先去数据库里\textbf{检索}相关段落 → 把段落喂给大模型（如ChatGPT）→ 大模型根据这些段落\textbf{生成}一个答案。\\[2pt]
\textbf{为什么需要RAG？}直接问ChatGPT，它不知道你论文里写了什么。但如果你先帮它找到相关段落，它就能基于论文内容回答了。\\[2pt]
\textbf{当前进度}：检索部分做好了，大模型API还没接。接上之后就是完整的RAG系统。
\end{termbox}

当前能做的：输入一个问题，返回论文中最相关的3段原文+页码+相关度。后续接上智谱/千问的API就能自动生成答案。

\begin{lstlisting}
# 使用方法（在命令行运行）
python 向量化存储.py --vectorize --ask "船舶燃气轮机课程知识体系如何构建？"
\end{lstlisting}


% ============================================================
\part*{第二部分：组会怎么讲}
% ============================================================

\section{开场（30秒）}

\begin{quote}
\kaishu
「这一周我主要完成了第一阶段的任务——从JSON数据的解析、清洗、分块、到向量化存储，全流程跑通了。另外我在基础功能之上做了三个优化：数据清洗（修复了OCR碎片问题）、混合检索（语义+关键词两路融合），以及搭好了RAG问答的框架。」
\end{quote}

% ============================================================
\section{数据概况（1分钟）}

\begin{quote}
\kaishu
「当前测试数据是一篇关于船舶燃气轮机课程知识体系的论文，4页，OCR标注后有246个文本块。

我在跑的过程中发现两个数据质量问题。

第一个是OCR标注把一些句子切断了，比如"构建方法"被切成了两半，只剩下"建方法"3个字，这种碎片会影响后面的检索质量。我在代码里加了自动合并——不足10个字的碎片自动拼到相邻块里，246块清洗到了228块。

第二个问题是双栏PDF的左右栏文字交叉，这个需要拿到标注框的x坐标才能解决，我想问一下郭老师JSON里是否有这个字段。」
\end{quote}

% ============================================================
\section{技术亮点（2分钟）}

\begin{quote}
\kaishu
「除了跑通基础流程，我额外做了三件事：

\textbf{第一，数据清洗。}把OCR碎片自动合并了，合并了18个碎片，质量检测从7个问题变成0个。清洗后分块更紧凑，chunk从21个减少到16个。

\textbf{第二，混合检索。}基础版只有语义检索（就是把文字转成数字来找相似段落），我加了BM25关键词检索（直接按关键词匹配），然后用RRF把两路结果融合排序。
举个例子：搜"布鲁姆认知过程"，语义检索找不太到，但关键词检索能直接命中。反过来搜"船舶动力装置"，关键词完全匹配不到，但语义检索能找到相关内容。

\textbf{第三，端到端问答框架。}现在可以输入一个问题，系统自动检索出最相关的段落和出处。后面接上大模型API就能自动写答案了，框架已经搭好。」
\end{quote}

% ============================================================
\section{下一步计划（30秒）}

\begin{quote}
\kaishu
「下一步有三件事：
\begin{enumerate}[nosep]
  \item 等全量的1000本手册JSON，把数据规模扩大
  \item 接入大模型API（智谱清言或者通义千问），实现真正的自动问答
  \item 开始调研Graph RAG，为后面的论文做准备
\end{enumerate}
」
\end{quote}

\begin{termbox}{什么是 Graph RAG？}
普通RAG只是把文字段落存起来搜索。Graph RAG在此基础上加入了\textbf{知识图谱}——把实体和关系用图的方式连起来（比如「BGE-m3 → 是一种 → 向量化模型」），这样能做更精准的推理。这是我们项目后续要做的、也是拿来发论文的方向。
\end{termbox}

% ============================================================
\section{刘超可能问的问题}

\subsection{Q: 为什么用BGE-m3不用其他模型？}
\begin{quote}
\kaishu
「BGE-m3是目前中文语义检索效果最好的开源模型之一，在MTEB中文排行榜上排名靠前。它支持多语言，参数量不大，普通电脑就能跑，不需要GPU服务器。」
\end{quote}

\begin{termbox}{MTEB 是什么？}
MTEB = Massive Text Embedding Benchmark，是评测向量化模型好不好的排行榜。类似于手机芯片的跑分排行。BGE-m3在中文任务上排名靠前。
\end{termbox}

\subsection{Q: 混合检索比纯语义检索好多少？}
\begin{quote}
\kaishu
「在精确关键词查询上明显更好。比如搜"布鲁姆认知过程"，纯语义只有0.387的相似度，排不到前面。但BM25通过关键词精确匹配到了，得分4.7分，直接排第一。两种方法互补，覆盖了各自的盲区。」
\end{quote}

\subsection{Q: 为什么分块用500字符？}
\begin{quote}
\kaishu
「这是经验值。太短信息不够，大模型理解不了上下文；太长噪音太多，搜出来一大段里面可能就一句话相关。500字符大约250个中文字，差不多一个自然段。后面数据量大了可以做实验调优。」
\end{quote}

\subsection{Q: RRF的70/30比例怎么定的？}
\begin{quote}
\kaishu
「0.7是业界常用的默认值。因为大多数情况下语义理解更重要，关键词只是补充。后续可以在测试集上调参。」
\end{quote}

\subsection{Q: 1000本手册能跑得动吗？}
\begin{quote}
\kaishu
「向量化是一次性的，1000本手册估计几万个chunk，可能需要几个小时跑完。但存好之后检索是毫秒级的。如果嫌慢可以用GPU加速。」
\end{quote}

\subsection{Q: 双栏文本交叉的问题解决了吗？}
\begin{quote}
\kaishu
「碎片合并的问题解决了，但双栏文字交叉的问题还没有完全解决。需要拿到标注框的x坐标，用x坐标判断这段文字属于左栏还是右栏，然后分别排序。我想确认一下JSON数据里是否有x坐标。」
\end{quote}

% ============================================================
\section{代码在哪}

\begin{lstlisting}
# GitLab 仓库（浏览器打开这个网址）
http://117.78.35.35/rag1/plant-graph-rag

# 你的分支（左上角下拉框切换到这个分支就能看到你的代码）
jiwenlong/vectorize-pipeline

# 你的代码文件
src/jiwenlong/pipeline.py

# 本地运行命令
cd 代码/
uv run python 源代码/纪文龙/向量化存储.py --vectorize --search
uv run python 源代码/纪文龙/向量化存储.py --vectorize --ask "你的问题"
\end{lstlisting}

\end{document}
